% =====================================================================
%  Temporal Self-Consistency: metric, measurements, and findings
%  Self-contained section; \input{} into the main paper.
%  Requires: booktabs, amsmath, multirow
% =====================================================================

\section{Measuring Temporal Self-Consistency}
\label{sec:te}

\subsection{Metric}
\label{sec:te-metric}

An agent is \emph{temporally self-consistent} if the action it takes at turn $t$ is one
it could already have predicted several turns earlier. Let $a_t$ be the action emitted
at turn $t$ of a trajectory $\tau$, tokenised as $a_t = (a_t^1,\dots,a_t^{n_t})$. We
truncate the dialogue at each of the $k$ preceding turns and ask the model to
\emph{forecast} the action at $t$; the forecast ensemble is the uniform mixture
\begin{equation}
  q^{F}(\cdot) \;=\; \frac{1}{k}\sum_{j=1}^{k}
  \pi_\theta\!\left(\cdot \;\middle|\; \tau_{<t-j}\right),
  \label{eq:qf}
\end{equation}
while $\pi^{0}(\cdot) = \pi_\theta(\cdot \mid \tau_{<t})$ is the model's
\emph{on-the-spot} distribution given the full preceding prefix. The per-token
temporal gain is
\begin{equation}
  g_t \;=\; \frac{1}{n_t}\Big[\log q^{F}(a_t) - \log \pi^{0}(a_t)\Big],
  \label{eq:gain}
\end{equation}
averaged over the target turns of $\tau$. Both terms are evaluated token-by-token over
the full vocabulary; $g_t = 0$ means the forecast made $k$ turns in advance is as
confident in the realised action as the model is once it is there, and $g_t < 0$
quantifies the missing foresight. We use $k=3$ and mixture weight $\eta = 0.5$.

\paragraph{The metric is a difference, and must be reported as one.}
Eq.~\eqref{eq:gain} rises under two mechanisms with opposite meaning: (A)
$\log q^{F}$ rises, i.e.\ the model genuinely anticipates the action earlier --- the
intended effect; or (B) $\log \pi^{0}$ falls, i.e.\ the model has moved away from
$a_t$ and no longer finds it likely even in context, so the gain rises with no
improvement in foresight. The two are empirically separable and, as
\S\ref{sec:te-findings} shows, they diverge in our data. We therefore report
$\Delta \log q^{F}$ and $-\Delta \log \pi^{0}$ alongside every $\Delta g$, and read a
gain unaccompanied by a rise in $\log q^{F}$ as an artefact.

\subsection{Evaluation protocol}
\label{sec:te-protocol}

Because $g_t$ is evaluated \emph{on} a trajectory, the choice of trajectories is part
of the measurement. Scoring each checkpoint on its own rollouts --- the quantity
logged during training --- is not viable: a model assigns near-certainty to its own
outputs (mean per-token $\log \pi^{0}$ of $-0.02$ against $-1.30$ on held-out
trajectories), so the denominator of Eq.~\eqref{eq:gain} collapses and $g_t$
degenerates into ``how well does the model predict text it wrote itself''. Holding
the \emph{weights fixed} and changing only the trajectory set moves the untrained base
model by $+3.53$, larger than the training effect of any checkpoint we evaluate.

All models are therefore scored on one frozen set of $60$ held-out trajectories
($1{,}428$ target turns), admitting only episodes that terminated with full task
reward, stratified to cover $20$ of the $26$ task types in the test split, and with
provenance rotated evenly across three independently trained runs ($21/20/19$) so that
no system is scored predominantly on its own outputs. Excluding the trajectories
contributed by the system under evaluation leaves all conclusions unchanged and
slightly strengthens the positive effects (CoPE at step $250$:
$+2.080 \rightarrow +2.211$; GRPO: $-0.630 \rightarrow -0.688$). Target turns within a
trajectory are correlated, so all intervals are $95\%$ bootstrap CIs over $4{,}000$
resamples at the trajectory level, computed on paired per-trajectory differences
against the base model.

\subsection{Results}
\label{sec:te-findings}

We evaluate an untrained base model and two systems at five checkpoints each: plain
GRPO, and CoPE, which augments the same RL objective with an auxiliary term training
the policy to predict its own next $k$ actions ($k=3$).

\begin{table}[t]
\centering
\small
\setlength{\tabcolsep}{5pt}
\begin{tabular}{llccc}
\toprule
System & Step & $\Delta g$ (95\% CI) & $\Delta \log q^{F}$ (95\% CI) & $-\Delta \log \pi^{0}$ \\
\midrule
\multirow{5}{*}{GRPO}
 & 50  & $+0.368$ \; $[+0.27, +0.47]$ & $+0.587$ \; $[+0.50, +0.69]$ & $-0.220$ \\
 & 100 & $+0.020$ \; $[-0.15, +0.19]$ & $+0.261$ \; $[+0.08, +0.44]$ & $-0.241$ \\
 & 150 & $-0.361$ \; $[-0.61, -0.13]$ & $-0.528$ \; $[-0.79, -0.26]$ & $+0.167$ \\
 & 200 & $-0.493$ \; $[-0.75, -0.24]$ & $-1.008$ \; $[-1.28, -0.73]$ & $+0.515$ \\
 & 250 & $-0.630$ \; $[-0.95, -0.35]$ & $-1.255$ \; $[-1.58, -0.97]$ & $+0.625$ \\
\midrule
\multirow{5}{*}{CoPE}
 & 50  & $+1.286$ \; $[+1.08, +1.50]$ & $\mathbf{+1.985}$ \; $[+1.75, +2.26]$ & $-0.700$ \\
 & 100 & $+1.836$ \; $[+1.59, +2.08]$ & $+1.819$ \; $[+1.58, +2.08]$ & $+0.018$ \\
 & 150 & $+1.946$ \; $[+1.65, +2.24]$ & $+1.461$ \; $[+1.21, +1.71]$ & $+0.486$ \\
 & 200 & $+2.085$ \; $[+1.83, +2.35]$ & $+1.501$ \; $[+1.27, +1.74]$ & $+0.584$ \\
 & 250 & $\mathbf{+2.080}$ \; $[+1.81, +2.34]$ & $+1.425$ \; $[+1.17, +1.71]$ & $+0.656$ \\
\bottomrule
\end{tabular}
\caption{Change in temporal gain relative to the untrained base model
($g = -1.74$, $\log q^{F} = -3.04$, $\log \pi^{0} = -1.30$), on $60$ held-out
successful trajectories ($1{,}428$ target turns). By construction
$\Delta g = \Delta \log q^{F} - \Delta \log \pi^{0}$. Bold marks the divergence
discussed in the text: CoPE's headline gain is largest at step $250$, but its
forecast term peaks at step $50$.}
\label{tab:te}
\end{table}

Table~\ref{tab:te} supports four observations.
\textbf{RL alone erodes temporal self-consistency.} Plain GRPO is briefly positive at
step $50$ ($+0.368$), loses significance by step $100$, and reverses to
$-0.630\,[-0.95,-0.35]$ at step $250$; the decline is driven by the forecast term
itself ($\Delta\log q^{F}\!: -0.528 \rightarrow -1.255$), so task reward makes the
policy \emph{less} able to anticipate its own actions even as its success rate rises.
\textbf{CoPE reverses this, via the intended mechanism.} CoPE is significant at every
checkpoint and reaches $+2.080\,[+1.81,+2.34]$, separating from GRPO by $2.7$ nats per
token with non-overlapping intervals throughout; at step $100$ the gain is almost
purely mechanism~(A) ($-\Delta\log\pi^{0} = +0.018$).
\textbf{Measured gain and underlying capability diverge later in training.} Beyond
step $50$ CoPE's forecast term declines monotonically ($+1.985 \rightarrow +1.425$)
while the headline gain rises, so by step $250$ roughly a third of the measured gain
reflects the policy moving away from the held-out actions rather than better
foresight --- foresight is acquired early and then plateaus. This is a limitation of
the metric rather than of the method, and the reason we treat the decomposition as
mandatory.
\textbf{The comparison is partly definitional.} CoPE's auxiliary objective trains the
model to predict its own next $k$ actions, close to what Eq.~\eqref{eq:gain} measures;
the result establishes that plain RL does not improve this quantity, but is weak
evidence of foresight beyond the form directly optimised. Whether it transfers to
behaviour is a question for the downstream task results.

\subsection{Limitations}
\label{sec:te-limitations}

$g_t$ is computed offline from stored evaluation trajectories with a single forward
pass. Against the quantity logged during training on the same checkpoints it tracks
sign, ordering and monotonicity but is systematically more negative by a factor of
$1.6$--$3.0$, because the training metric uses training-split rollouts and averages
over the $1.9$--$2.4$ forecast members available in a batch rather than the full
$k=3$; it is therefore a correlated proxy, not a reproduction. All conclusions are
relative comparisons on one frozen trajectory set and do not rely on numerical
agreement with the training-time quantity. We exclude turns whose action fails to
parse, but semantically invalid yet well-formed actions (e.g.\ opening a door that
does not exist) remain as targets, which inflates mechanism~(B); restricting to
successful trajectories mitigates but does not eliminate this. Finally, results come
from a single environment, base model, and training seed per system, and the bootstrap
intervals quantify sampling variation over trajectories, not across training runs.
